\section{The Unseen Hand}

The eye is a lazy witness. Show it what it expects to see, and it will sleep through a revolution. The ear is even worse. I have watched a student walk past a lecturing Synod clerk while humming a lullaby; the clerk did not look up. He was waiting for the punchline.

These spells do not fight. They \emph{nudge}. The Weave does not erase; it \emph{suggests}. And the observer, poor creature, accepts the suggestion because it is easier than doubt.

\subsection{On the Nature of Unseen Magic}

A spell of the unseen hand does not make a thing invisible. It makes it \emph{forgettable}. A guard who sees a veiled intruder will think \emph{``delivery boy''} and return to his dice. A merchant who finds a ledger page blank will think \emph{``clerk's error''} and close the drawer. The Weave does not lie. It merely omits.

The tags in this chapter are:

\begin{itemize}
    \item \textbf{[VEIL]} --- The caster blurs a single target's appearance. Not invisibility. \emph{Plausible denial}.
    \item \textbf{[SILENCE]} --- A zone or target becomes mute. The Weave does not steal sound; it \emph{postpones} it.
    \item \textbf{[ILLUSION]} --- The Weave paints a false image. It does not move unless the caster concentrates. It does not cast a shadow unless the caster remembers.
    \item \textbf{[SCRY]} --- The Weave lends the caster a distant eye. A glimpse of a room, a face, a document. Never the whole truth.
    \item \textbf{[MEMORY]} --- The Weave touches the past. A recollection may be blurred, sharpened, or---with care---removed. The Weave charges interest.
\end{itemize}

\subsection{Common Spells of the Unseen Hand}

\noindent
\begin{longtable}{p{0.2\textwidth} p{0.25\textwidth} p{0.45\textwidth}}
\caption{Spells of the Unseen Hand}\\
\midrule
\textbf{Spell} & \textbf{Tags} & \textbf{What It Does (in practice)}\\
\midrule
\endfirsthead
\multicolumn{3}{c}{\tablename\ \thetable{} -- continued}\\
\midrule
\textbf{Spell} & \textbf{Tags} & \textbf{What It Does (in practice)}\\
\midrule
\endhead
Hush & [SILENCE] & A zone the size of a dining table becomes mute. Shouts become whispers. Whispers become nothing. The Weave does not care what you say; it only cares that no one hears it. Lasts a few heartbeats. \\
Folded Face & [VEIL] & Your features blur. Not invisible---merely \emph{forgettable}. A guard will look at you and think \emph{``passing tradesman.''} A lover will feel a flicker of uncertainty. The effect ends if you speak your own name. \\
False Door & [VEIL][ILLUSION] & You paint an archway onto a stone wall. It looks real. It feels real to the touch, for a moment. It is not real. The Weave permits this deception, but it will remember that you lied. On a partial success, the door flickers when someone reaches for it. \\
The Unspoken Name & [SILENCE][TARGET] & You mute a single creature. They open their mouth. Nothing comes out. The Weave does not steal their voice; it merely \emph{misplaces} it for a few heartbeats. The target will find their voice again---usually at the worst possible moment. \\
Crowd's Shroud & [VEIL][AREA] & You blur the entire party. To any observer more than Near away, you are indistinguishable from market- goers, cargo crates, or stray dogs. The Weave does not make you invisible; it makes you \emph{boring}. \\
Distant Gaze & [SCRY] & You close your eyes and see a location you have visited before within a mile. The vision is grainy, like a reflection in rippled water. You cannot hear. You cannot act. Your body is helpless for the duration. \\
Gloss of the Forgotten & [MEMORY] & You touch a document, a ledger, or a letter. The Weave blurs a single word or number. The reader will see a smudge and assume the ink was wet. The Weave does not destroy evidence; it \emph{delays} discovery. \\
Phantom Footstep & [ILLUSION][AREA] & You create the sound of footsteps moving away from you---down a corridor, up a stair, into a crowd. A pursuer will follow the sound. For a few heartbeats, they will be looking in the wrong direction. \\
\hline
\end{longtable}

\subsection{On Failure, Unseen}

A veiling that fails is not a failure of invisibility. It is a failure of \emph{suggestion}. The witness sees what you tried to hide---and they know you tried to hide it.

\begin{tabular}{p{0.25\textwidth} p{0.7\textwidth}}
\textbf{Folded Face (Partial)} & Your features blur, but a distinctive scar or piece of jewelry remains visible. The witness may recognize the detail later. The GM gains a Story Beat. \\
\textbf{False Door (Miss)} & The archway appears, but it is conspicuously false---a painting, not a passage. Anyone who sees it will know magic was attempted. The GM gains two Story Beats. \\
\textbf{Crowd's Shroud (Partial)} & Your party blurs, but your footsteps remain audible. A guard tilts their head. They are now watching the area, not the crowd. \\
\textbf{Distant Gaze (Miss)} & You see nothing. Worse, the Weave shows you a \emph{wrong} image---a room you have never entered, a face you do not recognize. The GM may use this false vision as a complication later. \\
\textbf{Gloss of the Forgotten (Miss)} & The word you tried to blur becomes \emph{emphasized}. A subtle glow surrounds it. The reader will notice immediately. The GM gains a Story Beat. \\
\end{tabular}

\subsection{A Story from the Archivolt}

I was once hired to retrieve a document from a counting- house in Silkstrand. The document was a promissory note; the client was a merchant who had been cheated. I could have stolen it. I chose to erase it instead.

I cast [MEMORY] on the ledger page. The Weave obliged. The number vanished, replaced by a faint smudge that looked like old ink. The clerk closed the ledger. The debt was forgotten.

The Weave remembered, though. A month later, I received a letter. The letter contained a single word: \emph{``Interest.''} I have not opened a ledger since without leaving a pinch of salt on the spine.

\subsection{When to Use the Unseen Hand}

Use veiling when the truth would be inconvenient but a lie would be exhausting.

Use silence when the sound of your passage would be louder than your alibi.

Use illusion when you need a door that does not exist, or a guard who is not there.

Do not use memory lightly. The Weave is not a librarian. It does not stamp \emph{``withdrawn.''} It stamps \emph{``deferred.''} And the deferral always comes due.

\subsection{A Final Observation on Witness}

The unseen hand works best when there is a witness to deceive. A veiled face in an empty room is still a veiled face. A silenced argument between two people who are already alone is merely awkward.

But a veiled face in a crowd? A silenced footstep in a guarded hallway? The Weave thrives on witnesses. It uses their assumptions as fuel.

The best spell of the unseen hand is the one that convinces the witness that nothing happened at all.

\textit{--- Lysandra}
